% concatenated from main.tex
\documentclass[11pt]{article}

\usepackage[T1]{fontenc}
\usepackage{lmodern}
\usepackage{amsmath,amssymb,amsthm,mathtools}
\usepackage[margin=1in]{geometry}
\usepackage{microtype}
\usepackage{enumitem}
\usepackage{booktabs}
\usepackage[backend=biber,style=numeric,sorting=nyt,maxbibnames=20]{biblatex}
\addbibresource{references.bib}
\usepackage[hidelinks]{hyperref}

\hypersetup{
  pdftitle={The Subtractive Divisor Orbit: Unconditional Bounds, Parity Constraints, and a Conditional Framework},
  pdfauthor={Marco Mantovanelli}
}

\newtheorem{theorem}{Theorem}[section]
\newtheorem{lemma}[theorem]{Lemma}
\newtheorem{proposition}[theorem]{Proposition}
\newtheorem{corollary}[theorem]{Corollary}
\newtheorem{conjecture}[theorem]{Conjecture}
\newtheorem{hypothesis}[theorem]{Hypothesis}
\theoremstyle{definition}
\newtheorem{definition}[theorem]{Definition}
\theoremstyle{remark}
\newtheorem{remark}[theorem]{Remark}

\newcommand{\N}{\mathbb{N}}
\newcommand{\e}{\mathrm{e}}
\newcommand{\one}{\mathbf{1}}
\newcommand{\li}{\operatorname{li}}

\title{The Subtractive Divisor Orbit:\\
Unconditional Bounds, Parity Constraints, and a Conditional Framework}
\author{
Marco Mantovanelli\\
Independent Researcher\\
Germany\\
\texttt{marco@mantovanelli.de}
}
\date{}

\begin{document}
\maketitle

\begin{abstract}
Let $\tau(n)$ denote the number of positive divisors of $n$.  Starting
from an integer $x\geq 1$, consider the decreasing orbit
\[
 n_0=x,\qquad n_{j+1}=n_j-\tau(n_j),
\]
and let $a(x)$ be its hitting time of zero.  The average order of
$\tau$ suggests the order of magnitude $a(x)\asymp x/\log x$, but the
orbit is an endogenous sample of the divisor function.  We are not
aware of an unconditional estimate of this order.

We establish the exact displacement identity
$\sum_{j<a(x)}\tau(n_j)=x$ and the unconditional bounds
\[
 \frac{x}{(\log(2x))^3}\ll a(x)
 \leq \frac{3x}{8}+O\!\left(\frac{\sqrt{x}}{\log x}\right).
\]
The lower bound follows from the second moment of the divisor function.
The upper bound uses the fact that steps of size two occur only at
primes and that, apart from a finite exception, three such steps cannot
occur consecutively.  We also record a parity constraint: the orbit
changes parity exactly at square states.

On dyadic orbit segments we prove a local-to-global criterion, a
large-value truncation, and a quantitative implication from small
relative variance to a step-mass-saturating near-arithmetic progression.
We isolate the exact quantitative requirements in any Fourier
phase-rigidity argument, including the reduced modulus and the necessary
rate of phase concentration.  A conditional order-of-magnitude
reduction is then stated under two separate explicit hypotheses; neither hypothesis is
claimed here.  Finally, we provide reproducible numerical data through
$10^7$ and formulate several open problems concerning conditional
mixing, adaptive divisor ladders, and average behavior over starting
values.
\end{abstract}

\medskip
\noindent\textbf{Keywords:}
divisor function; arithmetic dynamics; deterministic recursion;
endogenous sampling; divisor moments; arithmetic progressions.

\medskip
\noindent\textbf{MSC2020:}
Primary 11A25; Secondary 11B37, 11N37, 11N60.

\section{Introduction}

For $n\in\N$, let
\[
 \tau(n)=\sum_{d\mid n}1
\]
be the divisor function.  We study the map
\[
 F(n)=n-\tau(n)
\]
and its forward orbit.  Thus, for a fixed starting value $x\in\N$,
\begin{equation}\label{eq:orbit}
 n_0(x)=x,\qquad n_{j+1}(x)=n_j(x)-\tau(n_j(x)).
\end{equation}
Since $1\leq\tau(n)\leq n$, the orbit remains non-negative, decreases
strictly while positive, and eventually reaches zero.  We write
\begin{equation}\label{eq:hitting-time}
 a(x)=\min\{j\geq 0:n_j(x)=0\}.
\end{equation}
This is sequence A155043 in the OEIS \cite{OEIS_A155043}.  To the best
of our knowledge, the unconditional bounds and the parity constraint
established below have not previously been recorded for this orbit.

The classical mean-value formula
\begin{equation}\label{eq:mean-tau}
 \sum_{n\leq X}\tau(n)=X\log X+(2\gamma-1)X+O(\sqrt X)
\end{equation}
suggests a mean step of logarithmic size and hence the order-of-magnitude
prediction
\begin{conjecture}\label{conj:main}
As $x\to\infty$ through the positive integers,
\[
 a(x)\asymp \frac{x}{\log x}.
\]
\end{conjecture}
Here and throughout, $f\asymp g$ means two-sided comparison by positive
absolute constants.  It does not mean $f/g\to 1$.

There are two immediate cautions.  First, $\log n$ describes the
\emph{average order} of $\tau(n)$; it is not its normal order.  Second,
the orbit does not sample integers independently.  Its current value of
$\tau$ determines the next sampling location.  Results for multiplicative
functions in intervals and arithmetic progressions
\cite{Shiu1980,MatomakiRadziwill2016,KlurmanMangerel2023,Mangerel2023,MatomakiRadziwillShaoTaoTeravainen2026}
therefore provide important comparison results but do not directly imply
Conjecture~\ref{conj:main}.
Likewise, correlation theorems for bounded multiplicative functions at
fixed shifts or almost all scales
\cite{MatomakiRadziwillTao2015,Tao2016,TaoTeravainen2019}
concern exogenously chosen samples rather than a path whose next state is
selected by its present divisor count.  Divisor concentration also
appears in quite different probabilistic settings
\cite{FordGreenKoukoulopoulos2023}; no direct transfer to the present
adaptive orbit is known.

\subsection{Unconditional result}

Two elementary comparison bounds are
\[
 \frac{x}{M(x)}\leq a(x)\leq\left\lceil\frac{x}{2}\right\rceil,
 \qquad
 M(x)=\max_{1\leq n\leq x}\tau(n).
\]
They follow by bounding every step from above by $M(x)$ and by using
$\tau(n)\geq2$ for $n\geq2$, respectively.  Their familiar coarse form
is $a(x)\gg_\varepsilon x^{1-\varepsilon}$ for each fixed
$\varepsilon>0$ and $a(x)\leq x/2+O(1)$.  Our main unconditional
theorem replaces both coarse estimates by explicit stronger bounds.

\begin{theorem}\label{thm:unconditional}
For every $x\geq 2$,
\begin{equation}\label{eq:unconditional}
 \frac{x^2}{\displaystyle\sum_{n\leq x}\tau(n)^2}
 \leq a(x)
 \leq \frac{3x}{8}+\frac{3}{8}\pi(\sqrt x)+\frac{11}{8}.
\end{equation}
Consequently,
\[
 \frac{x}{(\log(2x))^3}\ll a(x)
 \leq \frac{3x}{8}+O\!\left(\frac{\sqrt x}{\log x}\right).
\]
\end{theorem}

The lower bound uses only the exact displacement of the orbit and the
classical estimate
\begin{equation}\label{eq:second-moment}
 \sum_{n\leq X}\tau(n)^2\ll X(\log(2X))^3.
\end{equation}
The upper bound is a small structural gain: a step of size two occurs
exactly at a prime, and long runs of such steps are prohibited modulo
three.

The same parity phenomenon gives another unconditional constraint.

\begin{proposition}\label{prop:parity}
The parity of $n_j(x)$ changes exactly when $n_j(x)$ is a square.
Consequently, the orbit starting from $x$ changes parity at most
$\lfloor\sqrt x\rfloor$ times.
\end{proposition}

Thus uniform residue-class mixing, even modulo two, is not an automatic
or neutral model for this orbit.  Any future mixing statement must allow
for local arithmetic factors or condition on them explicitly.

\subsection{Structural scope}

We retain two useful parts of a structure-versus-randomness approach.
First, a sufficiently strong divisor-incidence regularity condition on a
dyadic orbit segment implies a logarithmic mean step.  Second, if the
step sizes on a segment have small relative variance, then almost all of
the displacement is carried by a dynamically compatible near-arithmetic
progression.  The latter conclusion holds with the segment mean $T$;
it does \emph{not} by itself imply $T\asymp\log N$.

We also explain precisely what a phase argument would have to prove.
For moduli growing with $N$, an $o(1)$ mean-square phase error is not
enough to force a residue class.  The error must be small relative to the
square of the reduced modulus.  Moreover, significant Fourier bias is
not the same as near-maximal bias.  These are open steps, not consequences
of the elementary identities.

\subsection{Organization}

Section~\ref{sec:exact} proves Theorem~\ref{thm:unconditional} and the
parity constraint.  Section~\ref{sec:dyadic} develops the dyadic
local-to-global criterion.  Section~\ref{sec:tails} gives the valid
large-value and single-bin reductions, with their limitations made
explicit.  Section~\ref{sec:ladders} proves the small-variance ladder
theorem.  Section~\ref{sec:phase} records quantitative phase lemmas.
Section~\ref{sec:conditional} gives a deliberately conditional two-hypothesis
reduction.  Numerical data and open problems appear in
Sections~\ref{sec:numerics} and~\ref{sec:open}.

\section{Exact dynamics and unconditional bounds}\label{sec:exact}

\subsection{Exact termination and displacement}

\begin{lemma}\label{lem:exact-energy}
For every $x\in\N$, the orbit in \eqref{eq:orbit} reaches zero and never
becomes negative.  Moreover,
\begin{equation}\label{eq:exact-energy}
 \sum_{j=0}^{a(x)-1}\tau(n_j(x))=x.
\end{equation}
More generally, if $0\leq r<s\leq a(x)$, then
\begin{equation}\label{eq:segment-energy}
 n_r(x)-n_s(x)=\sum_{j=r}^{s-1}\tau(n_j(x)).
\end{equation}
\end{lemma}

\begin{proof}
For $n\geq1$ one has $1\leq\tau(n)\leq n$.  Hence
$0\leq F(n)<n$, so iteration reaches zero after finitely many steps.
Summing $n_j-n_{j+1}=\tau(n_j)$ telescopes and gives both identities.
\end{proof}

It is useful to regard the left-hand side of \eqref{eq:exact-energy} as
the \emph{step mass} of the orbit.  We avoid the term additive energy,
which usually denotes a different, quadratic quantity.

\subsection{The second-moment lower bound}

\begin{proof}[Proof of the lower bound in Theorem~\ref{thm:unconditional}]
The positive values $n_0,\ldots,n_{a(x)-1}$ are distinct integers in
$[1,x]$.  By Lemma~\ref{lem:exact-energy} and Cauchy--Schwarz,
\[
 x^2
 =\left(\sum_{j<a(x)}\tau(n_j)\right)^2
 \leq a(x)\sum_{j<a(x)}\tau(n_j)^2
 \leq a(x)\sum_{n\leq x}\tau(n)^2.
\]
This proves the first inequality in \eqref{eq:unconditional}.  The
estimate \eqref{eq:second-moment}, a standard consequence of the
Dirichlet series $\sum\tau(n)^2n^{-s}=\zeta(s)^4/\zeta(2s)$, gives
$a(x)\gg x/(\log(2x))^3$; see, for example,
\cite{Apostol1976,Tenenbaum2015,JiaSankaranarayanan2014}.
\end{proof}

\subsection{A prime-run upper bound}

\begin{lemma}\label{lem:prime-runs}
Apart from the finite run $7,5,3$, a sequence of consecutive prime
states in the orbit has length at most two.
\end{lemma}

\begin{proof}
If $p$ is prime, then $\tau(p)=2$, so the next state is $p-2$.
Three consecutive odd prime states would therefore have the form
$p,p-2,p-4$.  One of these three integers is divisible by $3$; if all
are prime, that integer must equal $3$.  This gives only the run
$7,5,3$.
\end{proof}

\begin{proof}[Proof of the upper bound in Theorem~\ref{thm:unconditional}]
Partition the positive states of the orbit into four types.  Let $P$ be
the number of prime states, $S$ the number of states of the form $p^2$
with $p$ prime, $R$ the number of remaining states at least two, and
$U\in\{0,1\}$ indicate whether the state $1$ is visited.  Then
\[
 a(x)=P+R+S+U
\]
and, by the exact displacement identity,
\begin{equation}\label{eq:mass-types}
 x\geq 2P+4R+3S+U.
\end{equation}
Indeed, $\tau(n)=2$ exactly at primes, $\tau(n)=3$ exactly at squares
of primes, and every other $n\geq2$ has at least four divisors.

Prime states occur in runs separated by the $R+S$ composite states;
the possible terminal state $1$ cannot separate two prime runs.
Lemma~\ref{lem:prime-runs} therefore gives
\[
 P\leq 2(R+S)+3.
\]
Combining this with \eqref{eq:mass-types} gives
\begin{align*}
 3x-8a(x)
 &\geq -2P+S+4R-5U\\
 &\geq -3S-6-5U\\
 &\geq -3S-11.
\end{align*}
Therefore
\[
 a(x)\leq\frac{3x}{8}+\frac{3S}{8}+\frac{11}{8}.
\]
The orbit visits each integer at most once, so $S\leq\pi(\sqrt x)$.
This proves the exact upper bound.  The classical estimate
$\pi(y)\ll y/\log y$ gives its displayed
$O(\sqrt x/\log x)$ consequence.
\end{proof}

\subsection{Squares as parity gates}

\begin{proof}[Proof of Proposition~\ref{prop:parity}]
The classical pairing of divisors shows that $\tau(n)$ is odd if and
only if $n$ is a square.  Since
$n_{j+1}\equiv n_j-\tau(n_j)\pmod2$, the parity changes precisely at
square states.  There are at most $\lfloor\sqrt x\rfloor$ squares in
$[1,x]$, and the orbit is strictly decreasing.
\end{proof}

For later reference, the same observation has a Fourier formulation.
For any finite set $J$ of consecutive orbit indices,
\begin{equation}\label{eq:parity-phase}
 \frac1{|J|}\sum_{j\in J}
 \left|\e^{-\pi i\tau(n_j)}-1\right|^2
 =4\frac{\#\{j\in J:n_j\text{ is a square}\}}{|J|}.
\end{equation}
This deterministic identity is one reason that unconditioned residue
mixing is too crude a conjectural model.

\section{Dyadic orbit segments}\label{sec:dyadic}

\subsection{Hitting times and boundary errors}

For a fixed starting value $x$ and a real number $y\in[0,x]$, define
\[
 t_x(y)=\min\{j\geq0:n_j(x)\leq y\}.
\]
For $N>0$ with $2N\leq x$, the orbit block crossing $(N,2N]$ is
\[
 J_x(N)=\{t_x(2N),\ldots,t_x(N)-1\}.
\]
Write
\[
 V_x(N)=|J_x(N)|,
 \qquad
 E_x(N)=\sum_{j\in J_x(N)}\tau(n_j).
\]
Whenever $V_x(N)>0$, write in addition
\[
 T_x(N)=\frac{E_x(N)}{V_x(N)}
\]
for the mean step on the block.

\begin{lemma}\label{lem:hitting-boundary}
Uniformly for $1\leq y<x$,
\[
 n_{t_x(y)}=y+O(\sqrt y+1).
\]
Consequently, if $2N\leq x$, then
\begin{equation}\label{eq:dyadic-mass}
 E_x(N)=N+O(\sqrt N+1).
\end{equation}
\end{lemma}

\begin{proof}
Let $t=t_x(y)$.  If $t=0$, the assertion is immediate.  Otherwise put
$m=n_{t-1}>y$.  Since $n_t=m-\tau(m)\leq y$,
\[
 0<m-y\leq\tau(m)\leq2\sqrt m.
\]
For $m$ larger than an absolute constant this inequality implies
$m\leq2y+O(1)$, and hence $m-y=O(\sqrt y+1)$.  The same bound holds for
$y-n_t$.  Formula \eqref{eq:dyadic-mass} now follows by telescoping from
$t_x(2N)$ to $t_x(N)$.
\end{proof}

\subsection{A local-to-global criterion}

The natural dyadic scales for the orbit starting at $x$ are
\[
 N_k=\frac{x}{2^{k+1}}\qquad(k\geq0).
\]
These blocks partition the trajectory, up to their shared boundary
crossings and a final bounded range.

\begin{theorem}[Local-to-global criterion]\label{thm:local-global}
Suppose there are constants $0<c<C<\infty$ such that, for every
sufficiently large $x$ and every $k$ satisfying $N_k\geq\sqrt x$,
\begin{equation}\label{eq:local-mean-condition}
 c\log N_k\leq T_x(N_k)\leq C\log N_k.
\end{equation}
Then
\[
 a(x)\asymp\frac{x}{\log x}.
\]
\end{theorem}

\begin{proof}
For $N_k\geq\sqrt x$, Lemma~\ref{lem:hitting-boundary} gives
$E_x(N_k)\asymp N_k$.  Hence \eqref{eq:local-mean-condition} implies
\[
 V_x(N_k)\asymp\frac{N_k}{\log N_k}.
\]
The top block alone gives the lower bound $a(x)\gg x/\log x$.
For the upper bound, $\log N_k\geq\tfrac12\log x$ throughout the stated
range, so
\[
 \sum_{N_k\geq\sqrt x}V_x(N_k)
 \ll\frac1{\log x}\sum_{k\geq0}N_k
 \ll\frac{x}{\log x}.
\]
After the orbit first falls below $O(\sqrt x)$, at most $O(\sqrt x)$
further steps remain.  This is $o(x/\log x)$.
\end{proof}

The theorem is a sufficient local criterion.  No converse for each
individual dyadic block is asserted.

\subsection{A strong sufficient regularity condition}

For $d\geq1$, set
\[
 A_{x,N}(d)=\#\{j\in J_x(N):d\mid n_j(x)\}.
\]

\begin{definition}[Divisor-incidence regularity]\label{def:incidence}
The block $J_x(N)$ is divisor-incidence regular if
\begin{equation}\label{eq:incidence}
 \sum_{d\leq\sqrt{2N}}
 \left|A_{x,N}(d)-\frac{V_x(N)}d\right|
 =o\!\left(V_x(N)\log N\right)
\end{equation}
along the family of blocks under consideration.
\end{definition}

This is a deliberately strong sufficient condition, not a conjecture
that ignores the parity constraint of Proposition~\ref{prop:parity}.

\begin{proposition}\label{prop:incidence-mean}
If a family of dyadic blocks is divisor-incidence regular, then
\[
 T_x(N)\asymp\log N
\]
along that family.
\end{proposition}

\begin{proof}
For $n\in(N,2N]$,
\[
 \frac12\tau(n)
 \leq\sum_{\substack{d\leq\sqrt{2N}\\d\mid n}}1
 \leq\tau(n).
\]
Summing over $j\in J_x(N)$, interchanging the order of summation, and
using \eqref{eq:incidence}, we obtain
\[
 E_x(N)
 \asymp V_x(N)\sum_{d\leq\sqrt{2N}}\frac1d
      +o(V_x(N)\log N)
 \asymp V_x(N)\log N.
\]
Dividing by $V_x(N)$ proves the claim.
\end{proof}

\section{Large values and divisor scales}\label{sec:tails}

The second moment gives a useful truncation on every orbit block.  It
does not, by itself, identify a preferred divisor scale.

\begin{proposition}[Large-value truncation]\label{prop:tail}
Let $L\geq1$ and $2N\leq x$.  Then
\begin{equation}\label{eq:tail-general}
 \sum_{\substack{j\in J_x(N)\\ \tau(n_j)>L}}\tau(n_j)
 \ll \frac{N(\log(2N))^3}{L}.
\end{equation}
In particular, for every fixed $A>3$,
\[
 \sum_{\substack{j\in J_x(N)\\
                   \tau(n_j)>(\log N)^A}}
 \tau(n_j)=o(N)
\]
as $N\to\infty$, uniformly in the starting value $x$.
\end{proposition}

\begin{proof}
The orbit is strictly decreasing, so each integer is visited at most
once.  Since
$\tau(n)\one_{\tau(n)>L}\leq\tau(n)^2/L$, we have
\[
 \sum_{\substack{j\in J_x(N)\\ \tau(n_j)>L}}\tau(n_j)
 \leq\frac1L\sum_{N<n\leq2N}\tau(n)^2.
\]
Now apply \eqref{eq:second-moment} on the interval $(N,2N]$.
\end{proof}

The next proposition is the precise conclusion available from dyadic
pigeonholing.  It produces one non-negligible bin, not a bin carrying
almost all of the step mass.

\begin{proposition}[One-bin reduction]\label{prop:one-bin}
Fix $c>0$ and $A>3$.  Suppose $2N\leq x$ and
$B\subseteq J_x(N)$ satisfies
\[
 \sum_{j\in B}\tau(n_j)\geq cN.
\]
For all sufficiently large $N$, there are a dyadic number $T$ with
$1\leq T\leq(\log N)^A$ and a subset $B_T\subseteq B$ such that
\[
 T\leq\tau(n_j)<2T\quad(j\in B_T)
\]
and
\begin{equation}\label{eq:one-bin}
 T|B_T|\gg_{c,A}\frac{N}{\log\log N}.
\end{equation}
\end{proposition}

\begin{proof}
By Proposition~\ref{prop:tail}, values larger than $(\log N)^A$ carry
$o(N)$ step mass.  The remaining values of $\tau$ lie in
$O_A(\log\log N)$ dyadic ranges.  One such range carries
$\gg_{c,A}N/\log\log N$ mass.  On that range the mass is comparable to
$T|B_T|$.
\end{proof}

\begin{remark}[No critical scale from counting alone]\label{rem:no-critical}
If $T\leq\tau(n_j)<2T$ on a set of orbit indices, then the exact
displacement identity gives at most $O(N/T)$ such indices on a full
dyadic crossing.  Their total contribution can nevertheless be of
order $N$ for \emph{every} value of $T$.  Large steps are rarer but
heavier; small steps are lighter but may occur more often.  Thus the
counting bound does not imply $T\asymp\log N$.
\end{remark}

\section{Small variance and dynamic divisor ladders}\label{sec:ladders}

We now isolate a deterministic statement that does not require a
mixing hypothesis.  Throughout this section $J=J_x(N)$ is a non-empty
dyadic block, $V=|J|$, $E=E_x(N)$, and
\[
 T=\frac EV.
\]

\begin{definition}[Dynamic near-ladder]\label{def:dynamic-ladder}
Let $0<\eta<1$.  An ordered subsequence
\[
 m_1>m_2>\cdots>m_r
\]
of states from $J_x(N)$ is an $\eta$-\emph{dynamic near-ladder} with
parameter $T$ if
\begin{enumerate}[label=\textup{(\roman*)}]
 \item $|\tau(m_i)-T|\leq\eta T$ for every $i$;
 \item $m_i-m_{i+1}=\tau(m_i)+s_i$ with $s_i\geq0$;
 \item $\sum_{i<r}s_i\leq\eta E$.
\end{enumerate}
It is \emph{step-mass-saturating} if, in addition,
\[
 \sum_{i=1}^r\tau(m_i)\geq(1-\eta)E.
\]
\end{definition}

The slack $s_i$ is not arbitrary: it is exactly the step mass of orbit
states skipped between $m_i$ and $m_{i+1}$.  This makes the definition
dynamically compatible, unlike an arbitrary near-arithmetic
progression in $(N,2N]$.

\begin{theorem}[Small variance produces a dynamic ladder]
\label{thm:variance-ladder}
Assume $0<\varepsilon\leq1/16$ and
\begin{equation}\label{eq:variance-assumption}
 \frac1V\sum_{j\in J}(\tau(n_j)-T)^2
 \leq\varepsilon^2T^2.
\end{equation}
Then the states satisfying
\[
 |\tau(n_j)-T|\leq\sqrt\varepsilon\,T
\]
form a $3\sqrt\varepsilon$-dynamic near-ladder which is
step-mass-saturating.  Moreover, for all but
$O(\sqrt\varepsilon\,r)$ adjacent ladder pairs,
\begin{equation}\label{eq:ladder-gaps}
 m_i-m_{i+1}=T+O(\sqrt\varepsilon\,T).
\end{equation}
In particular,
\[
 r=(1+O(\varepsilon))V\asymp\frac ET.
\]
\end{theorem}

\begin{proof}
Let
\[
 \mathcal E=\{j\in J:|\tau(n_j)-T|>\sqrt\varepsilon\,T\}.
\]
Chebyshev's inequality and \eqref{eq:variance-assumption} give
\begin{equation}\label{eq:exception-count}
 |\mathcal E|\leq\varepsilon V.
\end{equation}
The step mass of this exceptional set is bounded by
\begin{align*}
 \sum_{j\in\mathcal E}\tau(n_j)
 &\leq T|\mathcal E|
 +|\mathcal E|^{1/2}
   \left(\sum_{j\in J}(\tau(n_j)-T)^2\right)^{1/2}\\
 &\leq (\varepsilon+\varepsilon^{3/2})TV
 \leq2\varepsilon E.
\end{align*}
Thus the complementary regular set carries at least $(1-2\varepsilon)E$
step mass.  List its states in orbit order as $m_1>\cdots>m_r$.
Equation \eqref{eq:exception-count} gives
$r=(1+O(\varepsilon))V$.

Between consecutive regular states, the orbit first makes the step
$\tau(m_i)$ and then possibly passes through exceptional states.  Hence
\[
 m_i-m_{i+1}=\tau(m_i)+s_i,
\]
where $s_i\geq0$ is the intervening exceptional step mass.  Summing in
$i$ gives
\[
 \sum_{i<r}s_i\leq2\varepsilon E.
\]
Because $\varepsilon\leq1/16$, all requirements of
Definition~\ref{def:dynamic-ladder} hold with
$\eta=3\sqrt\varepsilon$.

Finally, the number of indices with
$s_i>\sqrt\varepsilon T$ is at most
\[
 \frac{2\varepsilon E}{\sqrt\varepsilon T}
 =2\sqrt\varepsilon V
 =O(\sqrt\varepsilon r).
\]
For the remaining pairs, combine this with
$\tau(m_i)=T+O(\sqrt\varepsilon T)$ to obtain
\eqref{eq:ladder-gaps}.
\end{proof}

\begin{remark}
Theorem~\ref{thm:variance-ladder} is valid for the actual mean $T=E/V$.
It makes no assertion that $T\asymp\log N$.  Establishing that scale is
part of the original orbit problem, not a consequence of the ladder
construction.
\end{remark}

\section{Quantitative phase facts}\label{sec:phase}

Write $e(t)=\e^{2\pi i t}$.  For integers $q\geq2$ and
$h\not\equiv0\pmod q$, define
\[
 z_j=e\!\left(\frac{hn_j}{q}\right),
 \qquad
 u_j=e\!\left(-\frac{h\tau(n_j)}q\right).
\]

\begin{lemma}[Phase increment identity]\label{lem:phase-identity}
For every orbit index $j$,
\[
 z_{j+1}=z_ju_j,
 \qquad
 |z_{j+1}-z_j|=|u_j-1|.
\]
\end{lemma}

\begin{proof}
This is the exponential form of
$n_{j+1}=n_j-\tau(n_j)$ and the identity $|z_j|=1$.
\end{proof}

Near-maximal bias does force small phase increments.

\begin{proposition}[Near-maximal bias]\label{prop:max-bias}
Let $0\leq\delta\leq1$, and let $J=\{r,\ldots,s\}$ be a block of
$V=s-r+1$ consecutive orbit indices.  If
\[
 \left|\frac1V\sum_{j\in J}z_j\right|\geq1-\delta,
\]
then
\[
 \frac1V\sum_{j=r}^{s-1}|u_j-1|^2\leq8\delta.
\]
\end{proposition}

\begin{proof}
Let $\bar z=V^{-1}\sum_{j\in J}z_j$.  Since $|z_j|=1$,
\[
 \frac1V\sum_{j\in J}|z_j-\bar z|^2
 =1-|\bar z|^2\leq2\delta.
\]
The inequality
\[
 |z_{j+1}-z_j|^2
 \leq2|z_{j+1}-\bar z|^2+2|z_j-\bar z|^2
\]
and Lemma~\ref{lem:phase-identity} complete the proof.
\end{proof}

For growing moduli, the next quantitative form is essential.

\begin{proposition}[Phase concentration to a residue class]
\label{prop:phase-residue}
Let $Q=q/(h,q)$ and let $J$ be a non-empty finite orbit block of size
$V$.  Then
\begin{equation}\label{eq:phase-residue-quant}
 \frac1V\#\{j\in J:Q\nmid\tau(n_j)\}
 \leq\frac{Q^2}{16V}
 \sum_{j\in J}
 \left|e\!\left(-\frac{h\tau(n_j)}q\right)-1\right|^2.
\end{equation}
Consequently, along a family with $Q=Q(N)$, an exceptional proportion
$o(1)$ follows if the mean square on the right is $o(Q^{-2})$.
\end{proposition}

\begin{proof}
If $Q\nmid t$, then
\[
 \left|e\!\left(-\frac{ht}q\right)-1\right|
 \geq2\sin\frac\pi Q\geq\frac4Q.
\]
Summing this lower bound over the exceptional indices yields
\eqref{eq:phase-residue-quant}.
\end{proof}

The target phase $1$ detects the zero class modulo $Q$.  To detect an
arbitrary class $r\pmod Q$, the phase must instead be close to the
corresponding target $e(-hr/q)$.

\begin{lemma}[Correct CRT amplification]\label{lem:crt}
Let $Q_1,\ldots,Q_s$ be pairwise coprime and let
$Q=Q_1\cdots Q_s>M$.  If integers $t,t'\in[1,M]$ satisfy
\[
 t\equiv t'\pmod{Q_\ell}
 \qquad(1\leq\ell\leq s),
\]
then $t=t'$.
\end{lemma}

\begin{proof}
The Chinese remainder theorem gives $Q\mid(t-t')$, while
$|t-t'|<Q$.
\end{proof}

\begin{remark}[The remaining phase-rigidity gap]
Failure of divisor-incidence regularity gives a large \emph{sum} of
discrepancies.  It need not give one fixed relative discrepancy, and a
fixed relative Fourier bias need not be near-maximal.  Propositions
\ref{prop:max-bias} and~\ref{prop:phase-residue} therefore do not prove
that non-regularity produces small step variance.  In addition, when
several moduli are used, one needs the least common multiple of the
effective reduced moduli $q_\ell/(h_\ell,q_\ell)$ (their product in the
pairwise-coprime case), a common exceptional set, and a
summable total error.  These requirements are incorporated only as an
explicit hypothesis in the next section.
\end{remark}

\section{A conditional two-hypothesis framework}\label{sec:conditional}

The preceding results do not prove Conjecture~\ref{conj:main}.  This
section records a precise conditional reduction and, in particular,
keeps its two genuinely independent inputs separate.  On the non-empty
relative dyadic blocks below, define the
normalized divisor-incidence discrepancy
\[
 \mathcal D_x(N)=
 \frac{1}{V_x(N)\log N}
 \sum_{d\leq\sqrt{2N}}
 \left|A_{x,N}(d)-\frac{V_x(N)}d\right|.
\]
Only relative dyadic scales $N=N_k=x/2^{k+1}$ with
$N\geq\sqrt x$ will be used.

\begin{hypothesis}[Regularity-or-ladder dichotomy]\label{hyp:dichotomy}
There are functions
\[
 \rho,\eta:(1,\infty)\longrightarrow(0,1),
 \qquad \rho(N),\eta(N)\longrightarrow0,
\]
such that, for every sufficiently large $x$ and
every relative dyadic scale $N\geq\sqrt x$, at least one of the
following alternatives holds:
\begin{enumerate}[label=\textup{(\roman*)}]
 \item $\mathcal D_x(N)\leq\rho(N)$;
 \item $J_x(N)$ contains a step-mass-saturating
 $\eta(N)$-dynamic near-ladder with parameter $T_x(N)$.
\end{enumerate}
\end{hypothesis}

No size condition on $T_x(N)$ is included in alternative (ii).  In
particular, neither Proposition~\ref{prop:one-bin},
Theorem~\ref{thm:variance-ladder}, nor the phase lemmas supply the
logarithmic scale.

\begin{hypothesis}[Anti-ladder estimate]\label{hyp:antiladder}
With the same $\eta$ as in Hypothesis~\ref{hyp:dichotomy}, there is a
fixed $\delta\in(0,1)$ such that, for every sufficiently large $x$ and
every relative dyadic scale $N=N_k\geq\sqrt x$, every
$\eta(N)$-dynamic near-ladder with parameter $T_x(N)$
has
\[
 \sum_{i=1}^r\tau(m_i)\leq(1-\delta)E_x(N).
\]
\end{hypothesis}

Thus Hypothesis~\ref{hyp:dichotomy} says that non-regularity has a very
specific dynamically compatible consequence, at whatever mean scale
the orbit produces, whereas Hypothesis~\ref{hyp:antiladder} rules out
saturation by that consequence.  Neither is asserted unconditionally.

\begin{theorem}[Conditional order of magnitude]\label{thm:conditional}
If Hypotheses~\ref{hyp:dichotomy} and~\ref{hyp:antiladder} hold, then
\[
 a(x)\asymp\frac{x}{\log x}.
\]
\end{theorem}

\begin{proof}
Take $x$ large enough that $\eta(N)<\delta$ on every relative scale
$N\geq\sqrt x$.  If alternative (ii) of
Hypothesis~\ref{hyp:dichotomy} held, its ladder would carry at least
$(1-\eta(N))E_x(N)>(1-\delta)E_x(N)$, contradicting
Hypothesis~\ref{hyp:antiladder}.  Hence
$\mathcal D_x(N)\leq\rho(N)$ on every such block.

Because $\rho(N)\to0$, these blocks satisfy
Definition~\ref{def:incidence} uniformly.  The proof of
Proposition~\ref{prop:incidence-mean} then gives absolute constants
$0<c<C<\infty$ such that
\[
 c\log N\leq T_x(N)\leq C\log N
\]
on all the relevant scales.  Theorem~\ref{thm:local-global} completes
the proof.
\end{proof}

\begin{remark}
Theorem~\ref{thm:conditional} is a reduction, not evidence that either
hypothesis is true.  A successful phase-rigidity proof would have to
produce alternative (ii) with target residue classes, effective reduced
moduli, a common exceptional set, and errors strong enough for
Proposition~\ref{prop:phase-residue}.  A separate argument would then
have to establish the anti-ladder estimate.

The value of the reduction is therefore not that the two hypotheses
are presently accessible, but that it separates the dynamical rigidity
problem from the arithmetic anti-concentration problem.
\end{remark}

\section{Reproducible computations}\label{sec:numerics}

We computed $a(n)$ for every $1\leq n\leq10^7$ with a deterministic
C++17 program.  An Euler linear sieve first computes all divisor counts
in linear time, after which the recurrence
\[
 a(n)=1+a\bigl(n-\tau(n)\bigr)
\]
is evaluated in increasing order.  The program uses no randomness or
external data.  It checks the initial OEIS values, selected divisor
counts, every displayed checkpoint, the displacement identity
\eqref{eq:exact-energy}, and the equality between square visits and
parity changes.

Table~\ref{tab:normalizations} compares three normalizations.  Here
$\li(x)=\operatorname{Ei}(\log x)$, not the offset integral beginning
at $2$.  The middle normalization is merely a phenomenological
finite-range comparison; no limiting constant is asserted.

\begin{table}[htbp]
\centering
\small
\caption{Exact checkpoint values and rounded normalization ratios.}
\label{tab:normalizations}
\begin{tabular}{@{}r r c c c@{}}
\toprule
$x$ & $a(x)$
& $\dfrac{a(x)}{x/\log x}$
& $\dfrac{a(x)}{x/(\log x+\log\log x)}$
& $\dfrac{a(x)}{\li(x)}$\\
\midrule
$10$       & $3$      & $0.6908$ & $0.9410$ & $0.4866$\\
$10^2$     & $19$     & $0.8750$ & $1.1651$ & $0.6307$\\
$10^3$     & $116$    & $0.8013$ & $1.0255$ & $0.6531$\\
$10^4$     & $962$    & $0.8860$ & $1.0996$ & $0.7720$\\
$10^5$     & $7534$   & $0.8674$ & $1.0515$ & $0.7824$\\
$10^6$     & $65059$  & $0.8988$ & $1.0697$ & $0.8274$\\
$10^7$     & $556759$ & $0.8974$ & $1.0522$ & $0.8373$\\
\bottomrule
\end{tabular}
\end{table}

The orbit statistics in Table~\ref{tab:orbit-statistics} illustrate
two cautions.  First, the population standard deviation of the steps is
not small compared with their mean at these checkpoints, so the
hypothesis of Theorem~\ref{thm:variance-ladder} is not empirically
typical here.  Second, the proportion of even source states varies
strongly with the starting point, in agreement with the square-gate
constraint of Proposition~\ref{prop:parity}.
Here an even source is a positive pre-step state $n_j$ with
$n_j\equiv0\pmod2$; the terminal state zero is not counted.

\begin{table}[htbp]
\centering
\small
\caption{Statistics along four individual checkpoint orbits.}
\label{tab:orbit-statistics}
\begin{tabular}{@{}r r r r r r@{}}
\toprule
$x$ & mean $\tau$ & s.d. $\tau$ & even sources & squares & prime states\\
\midrule
$10^4$ & $10.395$ & $7.226$  & $90.23\%$ & $4$ & $16$\\
$10^5$ & $13.273$ & $10.186$ & $99.96\%$ & $2$ & $1$\\
$10^6$ & $15.371$ & $13.187$ & $97.32\%$ & $8$ & $273$\\
$10^7$ & $17.961$ & $16.768$ & $99.68\%$ & $8$ & $288$\\
\bottomrule
\end{tabular}
\end{table}

These data are compatible with the scale $x/\log x$ over the tested
range, but they do not distinguish an asymptotic law from slowly varying
corrections and do not prove Conjecture~\ref{conj:main}.  The source
program and the CSV files underlying the displayed tables accompany the
source package.  The full data through $10^7$ can be regenerated
deterministically.

\section{Open problems}\label{sec:open}

The unconditional and conditional results suggest several concrete
next steps.

\begin{enumerate}[label=\textup{\arabic*.},leftmargin=*]
 \item \emph{Close the unconditional gap.}  Improve the exponent three
 in the lower bound $a(x)\gg x/(\log x)^3$, or prove an upper bound
 $a(x)=o(x)$.  Either would require information beyond the global
 second moment and the elementary prime-run obstruction used here.

 \item \emph{Formulate parity-compatible regularity.}  Replace
 Definition~\ref{def:incidence} by a local-factor model that conditions
 on the parity phase, or on the most recent square crossing, and still
 yields $T_x(N)\asymp\log N$.

 \item \emph{Prove a quantitative rigidity implication.}  Identify
 hypotheses under which a sum of divisor-incidence discrepancies gives
 near-maximal phase concentration at enough target residues.  The error
 must track the effective moduli $q/(h,q)$ and survive a simultaneous
 CRT amplification.

 \item \emph{Control adaptive ladders.}  Establish a non-concentration
 theorem for dynamically selected near-ladders rather than for a fixed
 arithmetic progression chosen independently of the divisor values.
 This is the missing content of Hypothesis~\ref{hyp:antiladder}.

 \item \emph{Average over starting values.}  Determine the mean and
 variance of $a(x)$ for $x\leq X$, and ask whether a typical-start
 theorem of order $x/\log x$ is more accessible than a uniform theorem.
 The full functional graph of $n\mapsto n-\tau(n)$ may retain useful
 information that is invisible on a single orbit.
\end{enumerate}

\section*{Data and code availability}

The accompanying source package contains the complete C++17 program,
its build instructions, exact normalization checkpoints, orbit parity
statistics, and the values of $a(n)$ for $1\leq n\leq10^4$.  The
program regenerates all CSV files deterministically through $10^7$.

\section*{AI disclosure}

OpenAI's ChatGPT was used for language editing, bibliographic
organization, formal cross-checking, and \LaTeX{} assistance.  The
author independently verified all mathematical statements,
computations, and bibliographic claims and assumes full responsibility
for the manuscript.

\printbibliography

\end{document}
